\section{Low-fidelity labels and multi-fidelity training}\label{sec:multifidelity}

Multi-fidelity learning combines a few expensive, accurate samples with many cheap, approximate ones
\citep{kennedy2000predicting,peherstorfer2018survey}; for neural surrogates this is usually realised by
composite networks that correct a low-fidelity prediction with high-fidelity data
\citep{meng2020composite,lu2022multifidelity,howard2023multifidelity}. Here the two fidelities are two
\emph{label generators} for the same model, and the Ritz layer of \cref{sec:method} makes the cheap one natural.
Training experiments are in \cref{sec:experiments}.

\subsection{Two label fidelities}\label{sec:06-defs}

\paragraph{High fidelity ($\HF$).} The label is the adaptive order-3 solution of \cref{sec:data}. Frequencies and
figures of merit are taken from the order-3 label mesh; the training field target is the $M$-orthogonal
projection of the order-3 modes onto the order-2 space $\Vh$ of the model mesh.

\paragraph{Low fidelity ($\LF$).} The label is a single order-2 solve directly on the model mesh: no adaptivity and
no transfer between meshes. Because the label already lies in the model space, the field target carries no
projection error (the projection residual of the stored label is of order $10^{-7}$;
% src: docs/34 §1 (proj_err ~1e-7)
the $\HF$ projection residual, in contrast, is $0.10$--$0.42\,\%$ on the four geometries of \cref{tab:06-diff}).
% src: docs/34 §3 table (HF proj_err 0.10–0.42 %)

\paragraph{Why $\LF$ labels suit this model.} The Ritz layer computes eigenpairs of the model space itself, so the
frequencies it can attain are bounded below by the discrete order-2 eigenvalues $\lamh$ of that space
(\cref{prop:upper-bound}); these eigenvalues are the floor of the model space. $\LF$ frequencies are exactly these
eigenvalues, so the bound is tight at the label and a perfect basis reproduces it. An $\HF$ frequency is closer to
the exact eigenvalue and lies below the floor: the order-2 eigenvalue exceeds it by at most $0.033\,\%$ on the four
geometries of the first transfer study, and the $\LF$ frequency exceeds the $\HF$ one by $0.004$--$0.047\,\%$ on
identical meshes (positive in every measured case).
% src: docs/34 §4.1 table (FE floor ≤ 0.021/0.033/0.009/0.017 %); docs/34 §8.1 (+0.004–+0.047 %)
The $\HF$ field target is a different matter: it is the projection of the order-3 modes onto the model space,
and the Rayleigh quotients of these projected targets exceed the order-3 frequencies by $0.04$--$0.9\,\%$ (two
pillbox-with-pipes model meshes of 11\,300 and 13\,052 tetrahedra), an upper bound of the floor and a measure of the
quality of the projected target rather than of the floor itself.
% src: docs/30 §5 "Model-space frequency floor" (2.4e-3…8.6e-3, 3.7e-4…); EDITOR_NOTES 2

\subsection{The shared start mesh}\label{sec:06-mesh}

Mixing the two fidelities is meaningful only on the same mesh. The adaptive path starts from the curved
mesh with the first of a list of curvature-safety settings (the coarser $1.5$ first, the default $2.0$ as a fallback
when elements fold) and does not retry a finer setting after a mesh that exceeds the element limit; the $\LF$ path
uses the same rule.
% src: src/data_gen/field_labels.py start_mesh; docs/34 §8.1
Our first $\LF$ generator used the finer setting only and produced meshes up to $1.9$ times larger than the $\HF$
model meshes (\cref{sec:A06-meshes}); an earlier observation that the $\LF$ mesh ``does not shrink'' was a
misdiagnosis caused by this mismatch.
% src: docs/34 §8.1 (1.2–1.9×), §3 (corrected in §8.1)
With the shared rule the meshes are identical to the $\HF$ model meshes on 8 of 8 training and 6 of 6 test geometries,
the mean size over 40 $\LF$ samples drops from 9404 to 7639 tetrahedra ($-19\,\%$), and the cost ratio
$\HF$/$\LF$ rises from $2.3$ to $4.3$ (\SI{1242}{s} against \SI{291}{s} on the same eight meshes, two workers on a
shared 4-core machine under concurrent training load).
% src: docs/34 §8.1 (8/8, 6/6, 9404→7639 −19 %, 1242 s/291 s = 4.3, old ratio 2.3)
With the budget used here (at most 150\,000 label DOFs) the first order-3 solve exhausts the budget, so the
``adaptive'' $\HF$ label is a single order-3 solve on the start mesh; with a larger budget the $\HF$ model mesh may
be a refined level and the meshes would differ.
% src: docs/34 §8.1, §3, §5

\subsection{How the labels differ}\label{sec:06-diff}

On identical meshes, the $\LF$ frequency exceeds the $\HF$ one by $0.004$--$0.047\,\%$, and the part of an $\HF$
mode that falls outside the span of the six $\LF$ modes has an $M$-norm of $0.08$--$0.32\,\%$ for modes 0--3; the
model's own error is $10$--$50$ times larger than these differences.
% src: docs/34 §8.1 (freq +0.004–+0.047 %; field 0.08–0.32 %; 10–50×)
Quantities determined by the wall behave differently (\cref{tab:06-diff}). These values were measured on the earlier
meshes, where the $\LF$ and $\HF$ meshes coincided for only one of the four geometries, so they mix the effect of the
fidelity with that of the mesh; they were not repeated on identical meshes.
% src: docs/34 §3 (table; LF/HF tet counts differ for ids 5, 10, 13)

\begin{table}[t]
\centering
\caption{Relative difference of $\LF$ and $\HF$ labels (mode 0, four geometries), measured on the earlier meshes, which
coincided for only one geometry. Wall-dominated quantities differ much more than $\RoQ$.}
\label{tab:06-diff}
\small
\begin{tabular}{@{}lc@{}}
\toprule
Quantity & Range \\
\midrule
$\RoQ$ & $|\Delta|\le 0.04\,\%$ \\ % src: docs/34 §3 table (−0.005, +0.001, −0.039, +0.003 %)
$\Qz$ & $-0.76$ to $-1.37\,\%$ \\ % src: docs/34 §3 table
$\Epk/\Eacc$ & $-28$ to $+6.2\,\%$ \\ % src: docs/34 §3 table (+6.2, −28, −4.2, −13.6 %)
$\Bpk/\Eacc$ & $+5.4$ to $+8.1\,\%$ \\ % src: docs/34 §3 table
\bottomrule
\end{tabular}
\end{table}

Hence $\LF$ labels suffice for frequencies, volume fields and $\RoQ$ (differences well below the model error), whereas
$\Qz$ and the surface peaks need $\HF$ labels. The quantity-of-interest loss is disabled for the field model in this
study, so the experiments of \cref{sec:experiments} do not test wall quantities.
% src: docs/34 §3 (conclusion), qoi_weight = 0 for field3d

\subsection{Training mechanics}\label{sec:06-train}

Each sample carries a fidelity flag derived from its file metadata: a label produced with adaptivity, or with a label
order above the model order, is $\HF$; a single solve of the model order is $\LF$. Files of both kinds can be mixed.
% src: src/data/field_dataset.py label_fidelity
The span and frequency terms of the loss of \cref{sec:method} are weighted per sample,
\begin{equation}
  \mathcal{L}=\frac1B\sum_{b=1}^{B}\Bigl(w_{\varphi_b}\,\ell^{\mathrm{span}}_b
   +\lambda_f\,w_{\varphi_b}\,s_{\varphi_b}\,\ell^{\mathrm{freq}}_b\Bigr)+\cdots,
  \qquad \varphi_b\in\{\LF,\HF\},
\end{equation}
% src: src/training/lightning_module.py (fidelity_weights → w, fidelity_freq_scale → s; w_freq = w_fid * scale)
with fidelity weights $w_\LF,w_\HF$ and an extra scale $s_\LF,s_\HF$ on the frequency term alone, since $\HF$
frequencies lie below the floor of the model space; all default to one (no weighting). Per-fidelity errors are logged
separately.
% src: docs/34 §2
Three further mechanisms serve a scarce $\HF$ set: extra $\HF$ files that always enter the training split and never the
validation or test split; repetition of the $\HF$ geometries of the training split a given number of times per epoch; and
initialisation from the weights of a checkpoint, with fresh optimiser state and epoch counter, for fine-tuning an
$\LF$-trained model. A sweep of the fidelity scales and of the repetition count was not carried out, and the benefit of
mixed or fine-tuned training was not established at this scale (\cref{sec:experiments,sec:discussion}).
% src: docs/34 §2 (train_extra, hf_repeat, --init_from, frequency normalisation); §4.3, §6, §8.2 (sweep open; mixed and fine-tune runs inconclusive)
